\documentclass[11pt,a4paper]{article}
\usepackage{graphicx}
\usepackage{booktabs}
\usepackage{amsmath}
\usepackage{hyperref}
\usepackage{caption}
\usepackage{subcaption}
\usepackage{listings}
\usepackage{geometry}
\geometry{margin=1in}
\title{SVM Experiments with Hugging Face Model-Tree Relations}
\author{}
\date{}

\begin{document}
\maketitle

\begin{abstract}
Concise summary of the pipeline: export Hugging Face model-tree relations, extract features from cached DNA outputs, train a linear SVM probe to detect model-family signals, and evaluate with held-out test sets and cross-validation. This document summarizes implementation details, experimental design, train/test splitting strategy, and includes placeholder figures and tables for presentation.
\end{abstract}

\section{Overview}
The workflow has three main parts:
\begin{enumerate}
  \item Export model relations from Hub model-card metadata to form related-model groups.
  \item Extract feature vectors from precomputed DNA outputs for each model/run.
  \item Train and evaluate an SVM classifier that predicts membership in a related-model group.
\end{enumerate}

\section{Exporting the Hugging Face Model Tree}
Implementation summary:
\begin{itemize}
  \item For each model id in the input list, fetch repository metadata via the Hub API.
  \item Parse fields and card metadata for base-model relationships and config fields that indicate parent/base models.
  \item Add each model as a node, add an undirected edge for each discovered base relation (A--B).
  \item Collapse nodes into connected components to produce related-model groups.
\end{itemize}

Algorithmic detail (union--find streaming style):
\begin{lstlisting}[basicstyle=\ttfamily\small]
for model in input_models:
    info = hf_api.model_info(model)
    for base in extract_base_models(info):
        union_find.add(model)
        union_find.add(base)
        union_find.union(model, base)
components = collect_connected_components(union_find)
write_jsonl(groups)
\end{lstlisting}

Why union--find:
\begin{itemize}
  \item Supports streaming discovery of edges without constructing the full adjacency list first.
  \item Near-constant amortized cost per union/find.
  \item Low memory overhead and simple to implement.
\end{itemize}

Output:
\begin{itemize}
  \item JSONL file where each line is an object with a ``models'' array representing a related group.
  \item A summary JSON with counts (model\_count, node\_count, edge\_count, component\_count, skipped\_models).
\end{itemize}

\section{Feature Extraction and SVM Implementation}
Feature sources:
\begin{itemize}
  \item Per-model cached DNA outputs (e.g., response vectors, log-prob vectors, per-probe scores).
  \item Aggregated statistics across probes: mean response, variance, top-$k$ token agreement, entropy, and any hand-engineered distances to canonical outputs.
  \item Optionally, metadata features (model size bucket, family label, tokenizer type) if available.
\end{itemize}

Feature pipeline (recommended):
\begin{enumerate}
  \item Load cached outputs and align probes across models and runs to a common order.
  \item Compute numeric features per model or per (model, run) sample.
  \item Apply standard scaling (zero mean, unit variance) fit on training set; persist scaler.
  \item Optionally apply PCA to reduce dimensionality (retain e.g., 95\% variance).
\end{enumerate}

SVM details:
\begin{itemize}
  \item Use a linear SVM (e.g., scikit-learn's \texttt{LinearSVC} or \texttt{SGDClassifier} with hinge loss) for interpretability and speed on high-dimensional features.
  \item Regularization: tune penalty $C$ over a small grid (e.g., $[0.01,0.1,1,10]$).
  \item Use class weighting or balanced sampling if classes are imbalanced.
  \item For probabilistic outputs, wrap the SVM with Platt scaling (e.g., \texttt{CalibratedClassifierCV}) if needed for ROC/PR curves.
\end{itemize}

Minimal training loop (sklearn-style):
\begin{lstlisting}[basicstyle=\ttfamily\small]
from sklearn.preprocessing import StandardScaler
from sklearn.svm import LinearSVC
scaler = StandardScaler().fit(X_train)
Xtr = scaler.transform(X_train)
clf = LinearSVC(C=1.0, max_iter=10000).fit(Xtr, y_train)
\end{lstlisting}

\section{Experiment Outline}
Goals:
\begin{itemize}
  \item Evaluate whether DNA outputs contain signals that correlate with model-tree relations.
  \item Quantify detection performance (per-group and overall).
  \item Analyze robustness across sampling settings (temperature/top-p).
\end{itemize}

Experimental arms:
\begin{itemize}
  \item Baseline: random classifier or simple logistic regression.
  \item Linear SVM probe on raw features.
  \item Linear SVM + PCA dimensionality reduction.
  \item Cross-run generalization tests: train on one sampling regime, test on another.
\end{itemize}

Metrics:
\begin{itemize}
  \item Accuracy, precision, recall, F1 (micro and macro), ROC AUC, PR AUC.
  \item Confusion matrices aggregated at group-level.
\end{itemize}

\section{Train / Test Split Strategy}
Recommended approaches:
\begin{enumerate}
  \item \textbf{Group-aware holdout (preferred)}: Split by related-model group so that all models from some groups are only in train and other groups only in test. This evaluates generalization to unseen model families.
  \item \textbf{Model-level stratified split}: If labels are binary per-model (e.g., in-group vs out-group for a particular relation), stratify splits by label while ensuring models from the same group are not split across train/test to avoid leakage.
  \item \textbf{K-fold cross-validation with grouping}: Use GroupKFold where the group id is the connected-component id. This yields k folds where groups are exclusive per fold.
  \item \textbf{Per-run generalization}: Train on outputs from a subset of sampling settings (e.g., temperature=0.2) and test on others (temperature=0.3, different top-p).
\end{enumerate}

Practical splits:
\begin{itemize}
  \item Typical split sizes: 70\% train groups / 15\% val groups / 15\% test groups.
  \item For small group counts, prefer leave-one-group-out or repeated GroupKFold.
  \item When doing per-sample classification (many samples per model), ensure group-wise partitioning to prevent sample-level leakage.
\end{itemize}

\section{Experiment Details (Suggested Configuration)}
\begin{itemize}
  \item Input models: the full model list (N $\approx$ 11k). Optionally pre-filter by availability or size bucket.
  \item Feature vector dimensionality: depends on probes; example 512--4096. Use PCA if >1024.
  \item Classifier: LinearSVC with grid search on $C \in \{0.01,0.1,1,10\}$.
  \item Validation: GroupKFold with $k=5$ (groups = connected components).
  \item Repeat experiments with caching and with incremental checkpointing to measure wall time.
  \item Logging: store metrics per-fold, per-group summary, and top-k misclassified models for manual inspection.
\end{itemize}

\section{Figures and Tables (placeholders)}
\subsection{Figure suggestions}
\begin{figure}[ht]
  \centering
  \begin{subfigure}{0.48\textwidth}
    \centering
    \includegraphics[width=\linewidth]{figures/roc_curve_placeholder.pdf}
    \caption{ROC curve (SVM vs baseline)}
  \end{subfigure}
  \hfill
  \begin{subfigure}{0.48\textwidth}
    \centering
    \includegraphics[width=\linewidth]{figures/pr_curve_placeholder.pdf}
    \caption{Precision--Recall curve}
  \end{subfigure}
  \caption{Primary binary-classification performance plots.}
\end{figure}

\begin{figure}[ht]
  \centering
  \includegraphics[width=0.6\linewidth]{figures/stability_heatmap_placeholder.pdf}
  \caption{Heatmap showing probe stability or accuracy across sampling settings (temperature/top-p).}
\end{figure}

\subsection{Example tables}
\begin{table}[ht]
\centering
\caption{Aggregate classifier performance (mean $\pm$ std over folds)}
\begin{tabular}{lcccc}
\toprule
Model & Accuracy & Precision & Recall & F1 \\
\midrule
Baseline (random) & 0.50 & 0.50 & 0.50 & 0.50 \\
Logistic Regression & 0.72 $\pm$ 0.03 & 0.71 $\pm$ 0.04 & 0.70 $\pm$ 0.03 & 0.70 $\pm$ 0.03 \\
Linear SVM (C=1.0) & 0.78 $\pm$ 0.02 & 0.77 $\pm$ 0.03 & 0.76 $\pm$ 0.03 & 0.76 $\pm$ 0.02 \\
SVM + PCA (95\%) & 0.76 $\pm$ 0.03 & 0.75 $\pm$ 0.03 & 0.75 $\pm$ 0.04 & 0.75 $\pm$ 0.03 \\
\bottomrule
\end{tabular}
\end{table}

\begin{table}[ht]
\centering
\caption{Per-group breakdown (top 8 groups by size)}
\begin{tabular}{lrrrr}
\toprule
Group ID & \# Models & Accuracy & Precision & Recall \\
\midrule
group\_A & 42 & 0.92 & 0.90 & 0.95 \\
group\_B & 38 & 0.84 & 0.82 & 0.80 \\
group\_C & 30 & 0.79 & 0.75 & 0.78 \\
... & ... & ... & ... & ... \\
\bottomrule
\end{tabular}
\end{table}

\section{Practical Recommendations and Optimizations}
\begin{itemize}
  \item The dominant runtime cost is fetching model metadata; parallelize API calls with a thread pool and caching.
  \item Use group-aware splitting to avoid leakage from related models appearing in both train and test.
  \item Keep the classifier linear for interpretability and speed; examine coefficients to find discriminative features.
  \item Store intermediate caches and partial JSONL outputs to resume long runs.
  \item Save the scaler, PCA model, and trained SVM for reproducible evaluation pipelines.
\end{itemize}

\section{Appendix: Example Commands}
\begin{verbatim}
# Export model-tree relations (nohup recommended for long runs)
python export_hf_model_tree.py \
  --models-file huggingface_llm_list.jsonl \
  --output-file model_relations.jsonl \
  --summary-file model_relations_summary.json

# Train an SVM probe (sketch)
python scripts/train_svm_probe.py \
  --features-dir out/dna_features \
  --relations model_relations.jsonl \
  --output-dir out/svm_probe \
  --cv-folds 5 --svm-C 1.0
\end{verbatim}

\end{document}